% ===========================================================================
% TRUE-FIDELITY HARNESS
% Reproduces the book's real print environment so a prototype renders in the
% same type and the same margin a reader will actually see.
%   engine   : LuaLaTeX
%   fonts    : fontspec + mathptmx + newpxtext (Palatino) + helvet scaled 0.92
%   page     : 8x10in, outer 1.75in, marginparwidth 1.25in, marginparsep 0.25in
% All four lines copied from books/shared/tex/header-includes.tex.
% ===========================================================================
\usepackage{fontspec}
\usepackage{mathptmx}
\usepackage{newpxtext}
\usepackage[scaled=0.92]{helvet}
\usepackage[paperwidth=8in,paperheight=10in,
  top=0.5in,bottom=0.625in,inner=0.875in,outer=1.75in,
  headheight=14pt,headsep=18pt,includehead,footskip=24pt,
  marginparwidth=1.25in,marginparsep=0.25in,twoside]{geometry}
\usepackage{tikz}
\usetikzlibrary{positioning,calc,decorations.pathreplacing,decorations.pathmorphing}
\usepackage{xcolor}
\usepackage{ragged2e}
\pagestyle{empty}
\setlength{\parindent}{0pt}
\setlength{\parskip}{6pt}

\definecolor{crimson}{HTML}{A51C30}
\definecolor{ethzblue}{HTML}{1F407A}
\definecolor{amethyst}{HTML}{581C87}
\definecolor{emeraldpine}{HTML}{1A4D3E}

% Exactly the production rung: 2.6cm x 1.0cm, \rmfamily\scriptsize\bfseries,
% white 0.5pt separators, square corners, label inset 0.12cm.
% #1 y  #2 height  #3 label  #4 intensity  #5 accent
\newcommand{\rung}[5]{%
  \ifnum#4=0 \def\fillcol{black!5}\def\textcol{black!30}%
  \else \def\fillcol{#5!#4!white}%
    \ifnum#4>50 \def\textcol{white}\else \def\textcol{#5}\fi\fi
  \node[draw=white,line width=0.5pt,fill=\fillcol,minimum width=2.6cm,
        minimum height=#2 cm,anchor=south west,rounded corners=0pt,inner sep=0pt]
        (b) at (0,#1) {};
  \node[anchor=west,text=\textcol,font=\rmfamily\scriptsize\bfseries]
        at ($(b.west)+(0.12,0)$) {#3};
}
% Rung with a secondary right-hand line (cadence / unit of work).
\newcommand{\rungtwo}[6]{%
  \ifnum#5=0 \def\fillcol{black!5}\def\textcol{black!30}%
  \else \def\fillcol{#6!#5!white}%
    \ifnum#5>50 \def\textcol{white}\else \def\textcol{#6}\fi\fi
  \node[draw=white,line width=0.5pt,fill=\fillcol,minimum width=2.6cm,
        minimum height=#2 cm,anchor=south west,rounded corners=0pt,inner sep=0pt]
        (b) at (0,#1) {};
  \node[anchor=west,text=\textcol,font=\rmfamily\scriptsize\bfseries]
        at ($(b.west)+(0.12,0.11)$) {#3};
  \node[anchor=west,text=\textcol,font=\rmfamily\fontsize{5.4}{6}\selectfont,opacity=0.9]
        at ($(b.west)+(0.12,-0.14)$) {#4};
}
\newcommand{\spine}[5]{%
  \fill[#5!#4!white,rounded corners=1.5pt] (2.85,#1) rectangle (3.15,#2);
  \pgfmathsetmacro{\midy}{(#1+#2)/2}
  \ifnum#4>50 \def\sc{white}\else \def\sc{#5}\fi
  \node[rotate=90,text=\sc,font=\rmfamily\scriptsize\bfseries] at (3.0,\midy) {#3};
}

% Filler body prose so the margin object is judged against real type colour.
\newcommand{\bodypara}{%
Spatial memory is where an embodied system keeps what its sensors can no longer
see. A camera that observed a pallet two seconds ago has since looked elsewhere,
but the planner still needs to know the pallet is there, and needs to know how
much it should still trust that belief. The gap between those two questions is
the whole engineering problem. A transform published once and never expired will
be read as current forever, and a stale transform is not a degraded input but a
confident lie about where matter is.}
